% Reconstructed from the compiled lecture-note PDF.
\chapter{The singular locus is the circle}
\section{When the Hamiltonian forgets the control}

The control-dependent Hamiltonian is


\begin{displaytext}
Zu⟩ H2(X, ΛR; u) = −⟨ΛR, ⟨X, Zu⟩.
\end{displaytext}

\begin{displaytext}
Suppose this is independent of u on the whole triangle. The velocity set generated by UT spans the \
entire tangent plane to the determinant-one orbit at X. A covector that takes the same value on all \
these velocities must vanish. Hence \
ΛR = 0.
\end{displaytext}

\begin{displaytext}
Conversely, if ΛR = 0, the Hamiltonian has no explicit control dependence.
\end{displaytext}

\begin{statementbox}{Lemma}
Lemma 11.1. The Hamiltonian is independent of all u ∈UT if and only if ΛR = 0.
\end{statementbox}

\section{Determining the state on a singular interval}

\begin{displaytext}
Assume ΛR(t) = 0 on an open interval. Then also Λ′R(t) = 0. The Hamiltonian equation gives \
Λ1 −3 X = 0. 2λcostJ,
\end{displaytext}

\begin{displaytext}
The costate equation gives \
3 \
−Λ1 + 2λcostJ, X = 0. \
The first relation says the bracketed matrix is orthogonal to X; the second says it commutes with \
X. For a regular determinant-one traceless matrix, the centralizer is the line RX. The only matrix \
both parallel and orthogonal to X is zero because ⟨X, X⟩= −2. Therefore
\end{displaytext}

\begin{displaytext}
3 \
Λ1 = 2λcostJ.
\end{displaytext}

\begin{displaytext}
The nontriviality condition forces λcost ̸= 0, so the extremal is normal.
\end{displaytext}

\begin{displaytext}
Now use \
Λ′1 = [Λ1, X] = 0.
\end{displaytext}

\begin{displaytext}
Since Λ1 is a nonzero multiple of J, \
[J, X] = 0.
\end{displaytext}

The only X in the normalized orbit that commutes with J and satisfies ⟨J, X⟩< 0 is


X = J.


\begin{displaytext}
Then X′ = 0 forces the control matrix to commute with J. The unique control is the center of the \
triangle: \
1 1 1 1 \
u = 3, 3, 3 , Zu = 3J.
\end{displaytext}

\begin{displaytext}
Finally, \
g′ = gJ =⇒ g(t) = g0etJ. \
The curves g(t)s∗j are circular arcs, up to the fixed affine transformation g0.
\end{displaytext}

\begin{statementbox}{Definition}
Definition 11.2 (Singular locus). For normal scaling λcost = −1, the singular locus is
\end{statementbox}

\begin{displaytext}
Ssing = (g, X, Λ1, ΛR) : X = J, Λ1 = −3 ΛR = 0 . 2J,
\end{displaytext}

The group coordinate g ∈SL2(R) is arbitrary.


In upper-half-plane coordinates, X = J is the point


z = i.


\begin{statementbox}{Theorem}
Theorem 11.3 (Classification of full singular arcs). A Pontryagin extremal on which the Hamiltonian is independent of the entire control triangle is, up to affine transformation, a multicurve of circular arcs. Conversely, the circular multicurve gives such a singular extremal.
\end{statementbox}

\section{Why the circle is still a Pontryagin extremal}

The circle satisfies all first-order equations:


\begin{displaytext}
X = J, u = (1/3, 1/3, 1/3), ΛR = 0, Λ1 = −3J/2.
\end{displaytext}

This illustrates a central limitation of first-order necessary conditions. The circle is stationary under first-order variations but is not a local minimum. One needs a second-order variation.


\section{A second variation that lowers the area}

Consider the circular segment g0(t) = etJ.


Choose smooth compactly supported functions ψ1, ψ2 on an interval (t1, t2) and set


ψ1(t) ψ2(t) ! A(t) = . ψ2(t) −ψ1(t)


\begin{displaytext}
For small s, perturb by \
gs(t) = esA(t)etJ. \
Because A(t) is traceless, det esA(t) = 1, so gs(t) ∈SL2(R). The compact support keeps the endpoints \
fixed.
\end{displaytext}

The area cost can be written in a parameterization-invariant form using the Maurer–Cartan matrix g−1 dg: Z D cost(g) = −3 J, g−1g′E dt. 2 Expanding in s gives


Z t2 cost(gs) = cost(g0) −3s2 ψ′1ψ2 −ψ1ψ′2 dt + O(s3). t1


Choose nonnegative functions whose supports overlap and such that


ψ′1 > 0, ψ′2 < 0


on the overlap. Then ψ′1ψ2 −ψ1ψ′2 > 0,


so for sufficiently small s > 0, cost(gs) < cost(g0).


Because the circular curvature is strictly positive, sufficiently small C2 perturbations remain convex. Therefore the circular arc is not locally minimizing.


\begin{statementbox}{Theorem}
Theorem 11.4 (No singular arc in a global minimizer). A global Reinhardt minimizer cannot remain in Ssing on any interval of positive length.
\end{statementbox}

\begin{prooftext}
Proof. A singular interval is an affine image of a circular arc. The preceding compactly supported second variation lowers the area while preserving endpoint data and convexity, contradicting local optimality.
\end{prooftext}

\section{A single vertex segment cannot hit the singular locus}

Suppose an extremal uses one constant vertex control and reaches the singular locus at t = 0. The constant-control solution is analytic. At the singular point,


\begin{displaytext}
ΛR(0) = 0, Λ′R(0) = Λ′′R(0) = 0,
\end{displaytext}

but the third derivative is nonzero and points in a direction for which one of the other two vertex controls gives larger Hamiltonian. Hence the assumed vertex fails the maximum condition immediately. Therefore:


\begin{statementbox}{Lemma}
Lemma 11.5. An extremal segment with constant vertex control does not meet the singular locus.
\end{statementbox}

This does not rule out an infinite sequence of successively shorter vertex segments whose switching points converge to the singular locus. That is the chattering possibility.


\section{The trichotomy for a global minimizer}

A global minimizer is edge-extremal. We now know:


1. if it stays a positive distance from the singular locus, then it is finite bang-bang and hence a smoothed polygon; 2. it cannot remain on the singular locus for positive time; 3. no single constant-control segment reaches the singular locus.


Therefore the only unresolved behavior is an infinite switching sequence accumulating at the singular locus in finite time.


\begin{insightbox}{Proof roadmap}
\end{insightbox}

The rest of the proof is entirely devoted to excluding the following scenario:


\begin{displaytext}
t1 < t2 < · · · < tn ↗t∗, (X(tn), Λ1(tn), ΛR(tn)) −→ J, −3 0 . 2J,
\end{displaytext}

The control cycles among the vertices faster and faster while the state approaches the circular singularity.


\begin{insightbox}{Chapter summary}
\end{insightbox}

\begin{displaytext}
Full-triangle singularity forces ΛR = 0, X = J, and the central control u = (1/3, 1/3, 1/3); \
geometrically this is the circle. The circle satisfies Pontryagin’s first-order conditions but a \
compactly supported second variation lowers its area, so no minimizer contains a singular \
circular arc. The only remaining obstruction is chattering into the singular locus.
\end{displaytext}
